% Statement and cited-input packet. Target proof is separately embargoed.
% Frozen II source SHA256: 74497bfe327fae25c9cdcb17334b3bf42a9793a9df7c0a4404d97568f2f760f0
% Target owner: cdp:st:C.2.D4.3; tag: 01BY

\begin{proposition}[Rational connected triple covers have positive fibre Euler number]\hypertarget{tag-01BY}{}\label{cdp:st:C.2.D4.3}


Retain the torus-fibre, rational-normalization and specified
Hirzebruch-boundary morphism of \ref{cdp:st:C.2.RAT.1}. Suppose the horizontal
section $C$ is not contracted, its image $H$ is smooth rational,
the generic transversal along $H$ is an ordinary triple point,
and $C\to H$ has degree three. Then $e(S)\ge1$.
\end{proposition}

% Cited input cdp:st:C.2.D4.1; statement or labelled formula only.
\begin{proposition}[Global exclusion of connected positive-genus triple covers over a rational image]\hypertarget{tag-01BW}{}\label{cdp:st:C.2.D4.1}


Let $Y$ be a smooth compact complex threefold with
$b_2(Y;\Q)=0$, and let $f:Y\to B$ be a holomorphic map to a
smooth curve whose smooth nearby fibres are complex two-tori.
Fix a fibre $mS$, with $S$ reduced and irreducible, and restrict
the base to a small coordinate disc about it. Let
$H\subset S$ be a smooth rational curve along which the generic
transversal of $S$ is an ordinary triple point. Suppose that the
three generic normalization branches form one connected degree-three
cover $C\to H$, where $C$ denotes its smooth complete model.
Then $C$ is rational, and the cubic normal jet of $S$ vanishes
identically on at least one normal fibre over $H$.
In particular, a positive-genus connected degree-three cover is
impossible, and a primitive cubic normal jet is impossible for
any connected degree-three cover.
\end{proposition}

% Cited input cdp:st:C.2.RAT.1; statement or labelled formula only.
\begin{proposition}[Euler signs for two of the rational boundary classes]\hypertarget{tag-01CD}{}\label{cdp:st:C.2.RAT.1}


Let $f:Y\to B$ be a proper holomorphic map from a smooth compact
complex threefold with $b_2(Y)=0$, with connected fibres and smooth
general fibre a complex two-torus. Let $mS$ be a full fibre with
irreducible nonnormal reduced support. Write
$h=\eta\pi:Z\to T\to S$ for the minimal resolution of its
normalization and suppose that $Z$ is rational. Use the effective
divisor $K_Z+D=0$ of \SeriesRef{cdp:st:A.5.16}.

If a contraction to $\mathbf P^2$ sends $D$ to three concurrent
reduced lines, then $e(S)\ge1$. More precisely,
$e(S)=-1+N_0+e(\Gamma_{\rm red})\ge1+N_0$, with $N_0$ as in
(A.RT.12).

Consider also a specified morphism $\rho:Z\to\mathbb F_n$,
$n\ge1$, which is a sequence of point blowups in reverse and has
\[
 \rho_*D=2C_0+\sum_{j=1}^r a_jF_j,\qquad
 C_0^2=-n,\qquad
 a_j\ge1,\quad \sum_j a_j=n+2,                  \tag{C.RA.1}\hypertarget{tag-00WQ}{}\label{cdp:eq:C.RA.1}
\]
where the $F_j$ are distinct ruling fibres. Let $C$ be the
strict transform of $C_0$. Suppose $h(C)$ is a curve and the
generic transverse singularity there is a cusp $A_2$.
Then $r\le3$ and $e(S)\ge1$. If $h(C)$ is smooth, then
$r\le2$. These assertions allow singular normalization,
arbitrary vertical conductor multiplicities and a nonreduced
conductor-image scheme.
\end{proposition}

% Cited input cdp:st:C.2.RAT.3; statement or labelled formula only.
\begin{proposition}[Every uncontracted degree-one horizontal image has nonnegative Euler number]\hypertarget{tag-01CF}{}\label{cdp:st:C.2.RAT.3}


Retain the torus-fibre, rational-normalization and ruled-boundary
morphism of \ref{cdp:st:C.2.RAT.1}. Suppose that $C$ is not contracted and
that its actual map to the normalization of
$H=h(C)_{\rm red}$ has degree one. Let $r_{\rm dom}$ be the
number of original ruling fibres $F_j$ whose positive vertical
boundary contains a strict conductor component mapping onto
$H$. Then
\[
       0\le r_{\rm dom}\le2,\qquad
       e(S)\ge2-r_{\rm dom}\ge0.
\tag{C.RA.21}\hypertarget{tag-00X3}{}\label{cdp:eq:C.RA.21}
\]
There is no assumption on the generic transverse type, on the
local complete intersection property of $H$, or on the
nilpotent structure of the conductor image.
\end{proposition}

% Cited input cdp:st:C.2.RAT.5; statement or labelled formula only.
\begin{proposition}[A locally unibranch tangent-root component gives positive Euler number]\hypertarget{tag-01CH}{}\label{cdp:st:C.2.RAT.5}


Retain the torus-fibre, rational-normalization and specified
ruled-boundary morphism of \ref{cdp:st:C.2.RAT.1}. Suppose $C$ is not
contracted and maps with degree $d=2$ or $3$ onto a smooth
rational curve $H$ with generic ordinary-triple transversal.
Let $q_3\in H^0(H,\operatorname{Sym}^3N^*\otimes M)$ be
its cubic normal tensor, where $N=N_{H/Y}$ and
$M=\OO_Y(S)|_H$. Suppose $q_3$ is nonzero as a tensor at
every point of $H$. Its tangent-root divisor
$R_3\subset\PP(N)$ is then finite of degree three over $H$.
Let $R_d$ be the component corresponding to $C$ over the
generic point. If $R_d$ is locally unibranch at every point,
then
\[
                           e(S)\ge d-1>0.
\tag{C.RD.1}\hypertarget{tag-00XZ}{}\label{cdp:eq:C.RD.1}
\]
In particular this applies when $d=2$ and the only singularities
of $R_2$ are an $A_4$ cusp or two $A_2$ cusps.
\end{proposition}

